\section{Explicit constructions and obstructions}\label{sec:explicit}

\subsection{A virtual isomorphism with no invariant subgroup (F28)}\label{sec:f28}
\entry{F28}{S.~Sidki}
Must an isomorphism from an index-two subgroup of $F_2$ to a subgroup
of $F_2$ preserve a nontrivial subgroup of its domain? The following
example excludes even forward invariance, without a normality or
finite-generation assumption on the subgroup.

\begin{theorem}\label{prop:f28}
In $F=F(a,b)$ let $S=\langle a,bab^{-1},b^2\rangle$ and define
\[
 f(a)=b^{-2},\qquad f(bab^{-1})=b^{-1}a^2b,
 \qquad f(b^2)=b^{-1}a.
\]
Then $S$ and $f(S)$ have index two, $f$ is injective, and
$H\le S$, $f(H)\le H$ imply $H=1$.
\end{theorem}
\begin{proof}
The displayed generators of $S$ are the Schreier basis for the kernel
of $b$-parity, using transversal $\{1,b\}$. Represent $a,b$ in
$\mathrm{PSL}_2(\R)$ by
\[
 A=\begin{pmatrix}1&2\\0&1\end{pmatrix},\qquad
 B=\begin{pmatrix}1&0\\2&1\end{pmatrix}.
\]
Ping-pong on $\{|x|>1\}\cup\{\infty\}$ and $\{|x|<1\}$ shows
that this representation is faithful: every nonzero power of $A$
sends the second set into the first, and every nonzero power of $B$
sends the first into the second. Put
\[
 Q=\begin{pmatrix}0&-1\\2&1\end{pmatrix}.
\]
Direct multiplication gives
\[
 QAQ^{-1}=B^{-2},\quad Q(BAB^{-1})Q^{-1}=B^{-1}A^2B,
 \quad QB^2Q^{-1}=-B^{-1}A.
\]
The sign disappears projectively, so $f$ is injective. For the image
index, put $c=b^{-1}a$. The image is
$\langle b^2,c,bcb^{-1}\rangle$, an index-two Schreier subgroup
in the basis $b,c$.

If $h\in H$ has an integral determinant-one representative $M$, every
$M_n=Q^nMQ^{-n}$ is integral, since it represents $f^n(h)$ up to sign.
For $X=\left(\begin{smallmatrix}a&b\\c&d\end{smallmatrix}\right)$ define
\[
 E(X)=(8a-4b+2c-d)^2+7(4b+d)^2+7(2c-d)^2+49d^2.
\]
Expansion gives $E(QXQ^{-1})=E(X)$. The four linear forms determine
$X$, so the integral matrices of a fixed energy form a finite set.
Hence $M_i=M_j$ for some $i<j$, and $M$ commutes with $Q^{j-i}$.

Write $Q^n=u_nQ+v_nI$. From $Q^2=Q-2I$ one obtains
\[
 (u_0,v_0)=(0,1),\qquad (u_{n+1},v_{n+1})=(u_n+v_n,-2u_n).
\]
For $n\ge1$, $u_n$ is odd and $v_n$ is even. Thus $u_n\ne0$,
and commuting with a positive power of $Q$ implies commuting with $Q$.
The latter equation gives $c=-2b$, $d=a-b$, whence
\[
 1=\det M=a^2-ab+2b^2=(a-b/2)^2+7b^2/4.
\]
Integrality forces $b=c=0$, $a=d=\pm1$. Thus $h=1$ projectively
and, by faithfulness, in $F$.
\end{proof}

Related rational conjugations appear in work on self-similar
groups~\cite{kapovich-selfsimilar}. Simplicity there excludes
normal invariant subgroups; Proposition~\ref{prop:f28} concerns every
subgroup. The integral-matrix implication has a Lean proof, while the
free-group interpretation remains a written argument.
\evidence{F28}{proof.md}

\subsection{Short presentations and long Dehn output (H4)}\label{sec:h4}
\entry{H4}{A.~Miasnikov}
Can a finite presentation of a hyperbolic group be converted to a Dehn
presentation in polynomial time? We use explicit words in generators
and their inverses, and strict length-decreasing Dehn rules.

\begin{theorem}\label{prop:h4}
No uniform polynomial-time algorithm performs this conversion on all
finite hyperbolic-group presentations. This remains so when the inputs
are infinite non-elementary virtually free groups and the output may
use different generators.
\end{theorem}
\begin{proof}
Let a Dehn presentation have $r$ generators and longest relator length
$L\ge1$. A classical observation bounds the number $T(G)$ of torsion
conjugacy classes, including the identity, by
\begin{equation}\label{eq:torsion-dehn}
 T(G)\le\sum_{j=0}^{L-1}(2r)^j<(2r)^L.
\end{equation}
Indeed, take a shortest word $w$ in a nontrivial torsion conjugacy class.
It is cyclically reduced. If $|w|\ge L$, a Dehn-shortenable subword
of an identity power $w^e$ has length at most $L$ and lies in a cyclic
rotation of $w$. Shortening it contradicts conjugacy minimality. This
is the usual torsion-conjugacy argument~\cite[Theorem 3.27]{batty-papasoglu}.
The same reasoning applies to an explicit finite list of shortening
rules, using the largest left-hand-side length.

For $n\ge1$ take generators $x,t_0,\ldots,t_n$ and relators
\begin{equation}\label{eq:short-presentation}
 t_i=t_{i-1}^2\ (1\le i\le n),\qquad t_n=1,
 \qquad t_0xt_0^{-1}=x^2.
\end{equation}
Put $Q_n=2^n$, $M_n=2^{Q_n}-1$. Eliminating the $t_i$ gives
\[
 K_n\cong C_{M_n}\rtimes C_{Q_n},\qquad t_0xt_0^{-1}=x^2.
\]
To check equality rather than just a quotient, the relations give
$t_0^{Q_n}=1$ and $x^{M_n}=1$ and collect every word to
$x^it_0^j$, so the order is at most $M_nQ_n$. The indicated
semidirect product realizes that order and the presentation. Multiplication
by two has order $Q_n$ modulo $M_n$, since $0<2^j-1<M_n$ for
$0<j<Q_n$.

The conjugacy classes in $\langle x\rangle$ are the orbits of this
cyclic action, so there are at least
\[
 M_n/Q_n>2^{2^n-n-1}
\]
of them. They remain distinct in $H_n=K_n*\langle z\rangle$ by
free-product normal forms. The kernel of the retraction $H_n\to K_n$
has free basis $\{kzk^{-1}:k\in K_n\}$ and finite index $|K_n|$.
Its rank is at least six. Hence $H_n$ is infinite, non-elementary
and virtually free, in particular hyperbolic.

Adding $z$ to~\eqref{eq:short-presentation} gives explicit
generator-letter size $4n+9$, or $O(n\log n)$ bits with indexed
generators. Any explicit Dehn output of total size $m_n$ has at most
$m_n$ generators and longest rule length at most $m_n$. Therefore
\[
 m_n\log_2(2m_n)>2^n-n-1,
 \qquad m_n>\frac{2^n-n-1}{n+1}\quad(n\ge2).
\]
The second inequality follows by separating $m_n\le2^n$ and
$m_n>2^n$. Merely writing the output takes superpolynomial time
in the input length, regardless of its chosen generators.
\end{proof}

The torsion hypothesis and explicit-output convention matter. This
argument makes no claim about torsion-free input or compressed output.
The original finite-input argument, the classical torsion lemma and
related earlier output-size observations are separately credited in
the \record{problems/H4/audit.md}{audit}.
\evidence{H4}{infinite-input-proof.md}

\subsection{Primitive elements in free metabelian groups (M0)}\label{sec:m0}
\entry{M0}{E.~I.~Timoshenko and V.~Shpilrain}
The entry asks the analogue of F3 for $M_r=F_r/F_r''$: does preservation
of all primitive elements force an endomorphism to be an automorphism?
Primitive means membership in a free basis of the metabelian variety.

\begin{theorem}\label{prop:m0}
For every finite rank $r$, an endomorphism of $M_r$ preserving primitive
elements is an automorphism. It suffices to preserve the images of
primitive words of $F_r$.
\end{theorem}
\begin{proof}
Use the Laurent ring $R=\Z[X_1^{\pm1},\ldots,X_r^{\pm1}]$ and
abelianized left Fox derivative columns $D(w)$. Write $J_\phi$ for
the matrix with columns $D(\phi(x_j))$ and
$\lambda=(X_1-1,\ldots,X_r-1)$. The standard identities are
\[
 D(uv)=D(u)+\bar uD(v),\quad \lambda D(w)=\bar w-1,
 \quad D(\phi(w))=J_\phi\bar\phi(D(w)).
\]
Bachmuth's inverse-function criterion, in this convention, says that
$\phi$ is invertible exactly when $J_\phi$ is invertible over $R$
\cite[Lemma 1]{ggr}. A primitive word's derivative cannot specialize
to zero over a field, since it is a column of an invertible basis
Jacobian.

The induced integer matrix on abelianization preserves primitive
vectors and is therefore unimodular. Otherwise reduction modulo a
prime gives a nonzero kernel vector with a primitive integral lift,
obtained by lifting elementary matrices in $\SL_r(\F_p)$; its image
is divisible by $p$. Compose with a lifted inverse basis change to
assume $\phi$ is IA. Then
\begin{equation}\label{eq:ia-jacobian}
 \lambda J_\phi=\lambda,\qquad J_\phi(1,\ldots,1)=I.
\end{equation}
Ranks zero and one are immediate, so assume $r\ge2$.

If $\Delta=\det J_\phi$ is not a Laurent unit, augmentation one
implies it has at least two monomials. Choose a prime preserving two
nonzero coefficients and a maximal ideal containing its reduction.
Its residue field $K$ is finite. The images $z_i$ of $X_i$ are
nonzero, generate $K$ as a field, and are not all one. The associated
character has a nontrivial finite cyclic image. A Nielsen basis change
puts it in the form
\[
 \chi(x_1)=\cdots=\chi(x_{r-1})=1,\qquad
 \chi(x_r)=t\ne1,\qquad K=\F_p(t).
\]
Indeed, apply integer elementary operations to an exponent row for
a generator of that cyclic image; its gcd remains coprime to the
image order. The Fox basis change conjugates $J_\phi(\chi)$,
because $\phi$ is IA, so singularity is preserved. Equation~\eqref{eq:ia-jacobian}
now puts its kernel in $W=\Span_K(e_1,\ldots,e_{r-1})$.

For distinct $i,j<r$ and $P(T)=\sum a_kT^k\in\Z[T]$, fix the
other generators and set
\begin{equation}\label{eq:m0-transvection}
 \theta(x_i)=x_i\prod_{k\ge0}x_r^kx_j^{a_k}x_r^{-k}.
\end{equation}
This is an actual free-group automorphism preserving $\chi$, and
$J_\theta(\chi)=I+P(t)E_{ji}$. Since every element of $K$ is
$P(t)$, for $r\ge3$ these maps realize $\SL_{r-1}(K)$ on $W$.
Its transitivity on nonzero vectors shows that every nonzero kernel
vector is $D(w)(\chi)$ for a primitive word $w$. If $r=2$, the
nonzero kernel is the line $Ke_1$, and $w=x_1$ suffices.
The chain rule gives $D(\phi(w))(\chi)=0$, contradicting primitivity.
Thus $\Delta$ is a unit and the inverse-function criterion applies.
\end{proof}

The low-rank answers are prior; the all-rank finite-field detection is
the proposed additional step. The proof neither assumes that every
metabelian automorphism is tame nor substitutes preservation of whole
primitive systems for preservation of individual elements. Exact
finite-field witnesses include actual free bases and inverses. A later
constructor removes a search cutoff by an explicit Kronecker bound;
it gives termination without a useful complexity estimate.
\evidence{M0}{proof.md}
